% ============================================================
% CAPÍTULO 1: INTRODUCCIÓN
% ============================================================

\section{Motivación}
% Explica el problema que motiva tu TFG:
% - La dificultad de aprender programación
% - El riesgo de que los LLMs fomenten el "copypasteo"
% - La necesidad de un enfoque socrático
% TODO: Redactar


\section{Objetivos}

\subsection{Objetivo General}
% Describe el objetivo principal del TFG en un párrafo.
% Ejemplo: Desarrollar una plataforma web que integre un tutor de programación
% basado en IA con metodología socrática.
% TODO: Redactar

\subsection{Objetivos Específicos}
% Lista los objetivos concretos y medibles.
\begin{itemize}
    \item % TODO: Objetivo específico 1 (ej: Diseñar e implementar el backend y frontend)
    \item % TODO: Objetivo específico 2 (ej: Integrar una API de LLM)
    \item % TODO: Objetivo específico 3 (ej: Desarrollar el system prompt socrático)
    \item % TODO: Objetivo específico 4 (ej: Validar con usuarios reales)
\end{itemize}


\section{Estructura de la Memoria}
% Describe brevemente el contenido de cada capítulo de la memoria.
% Ejemplo: "Tras esta introducción, el Capítulo~\ref{CAP:ESTADO_ARTE} revisa
% el estado del arte..."
% TODO: Redactar (ajustar según los capítulos definitivos)
